\noindent\textbf{关系图阅读方法。}\quad 图中的表卡显示物理表和字段；以下\emph{Key Path} 按“子表外键 $\rightarrow$ 父表主键”书写。
\emph{Nullable} 描述子表外键是否可空，\emph{ON DELETE} 是现有快照的删除动作，均不等同于业务操作权限。
\emph{UNIQUE}、状态互斥、金额上限与时段不重叠另由唯一约束、事务或应用规则执行。

\begin{table}[!htbp]\centering\small
\caption{关系图与关系说明的阅读层次}\label{tab:erd-reading-levels}
\begin{tabular}{@{}L{30mm}L{50mm}L{69mm}@{}}\toprule
层次 / View & 回答的问题 & 阅读范围\\\midrule
全景 / Panorama & 有哪些业务簇和跨域依赖？ & 61 张表及跨域连线；用于总览，不在 A4 上逐字段阅读。\\
分域 / Domain ERD & 本域的表怎样相连？ & 每张分域图显示字段、PK/FK 及关系箭头，九域逐张放大检查。\\
关系说明 / Key Path & 外键方向和可空性是什么？ & 下列九张关系表列出现有约束及其业务语义；未选的完整关系见附属清单。\\
字段字典 / Data Dictionary & 每个字段的类型和默认值是什么？ & 独立附录逐字段列出类型、键、可空、默认值和索引。\\
\bottomrule\end{tabular}\end{table}

\subsubsection{组织、账户与授权：关键关系与约束}
组织边界与岗位授权分开保存。组织、诊所是权限范围的定位点，岗位角色通过中间表授予用户，操作人不能由患者资料代替。

\begingroup\footnotesize\setlength{\tabcolsep}{3pt}
\begin{longtable}{@{}L{43mm}L{32mm}L{24mm}L{49mm}@{}}
\caption{组织、账户与授权关系路径（Key Paths）}\label{tab:domain-key-01}\\\toprule
Child table / FK & Parent table / PK & Null; Delete & Business meaning / 业务含义\\\midrule\endfirsthead
\toprule Child table / FK & Parent table / PK & Null; Delete & Business meaning / 业务含义\\\midrule\endhead
\texttt{clinics}\newline\texttt{organization\_id} & \texttt{organizations}\newline\texttt{id} & NO\newline\texttt{RESTRICT} & 诊所归属组织，跨诊所授权须先确定组织范围。\\
\rowcolor{Pale}\texttt{role\_permissions}\newline\texttt{role\_id} & \texttt{roles}\newline\texttt{id} & NO\newline\texttt{CASCADE} & 角色与权限采用关联表，避免在角色记录中存放权限列表。\\
\texttt{role\_permissions}\newline\texttt{permission\_id} & \texttt{permissions}\newline\texttt{id} & NO\newline\texttt{CASCADE} & 同一权限可由多个角色复用。\\
\rowcolor{Pale}\texttt{user\_role\_assignments}\newline\texttt{user\_id} & \texttt{app\_users}\newline\texttt{id} & NO\newline\texttt{CASCADE} & 一个账户可持有多个角色授权，授权本身可独立留痕。\\
\texttt{user\_role\_assignments}\newline\texttt{clinic\_id} & \texttt{clinics}\newline\texttt{id} & YES\newline\texttt{RESTRICT} & 诊所范围可选；不应把空值自动解释为全网权限。\\
\rowcolor{Pale}\texttt{audit\_events}\newline\texttt{actor\_user\_id} & \texttt{app\_users}\newline\texttt{id} & YES\newline\texttt{SET NULL} & 审计事实与操作账户关联；账户停用不能抹去历史事件。\\
\bottomrule\end{longtable}\endgroup


\subsubsection{身份认证与通知：关键关系与约束}
身份凭证、会话和患者本人绑定采用不同生命周期；通知记录与业务预约关联，但通知成功不等于预约已确认。

\begingroup\footnotesize\setlength{\tabcolsep}{3pt}
\begin{longtable}{@{}L{43mm}L{32mm}L{24mm}L{49mm}@{}}
\caption{身份认证与通知关系路径（Key Paths）}\label{tab:domain-key-02}\\\toprule
Child table / FK & Parent table / PK & Null; Delete & Business meaning / 业务含义\\\midrule\endfirsthead
\toprule Child table / FK & Parent table / PK & Null; Delete & Business meaning / 业务含义\\\midrule\endhead
\texttt{auth\_identities}\newline\texttt{user\_id} & \texttt{app\_users}\newline\texttt{id} & NO\newline\texttt{CASCADE} & 第三方身份映射归属内部账户，外部标识不能充当业务主键。\\
\rowcolor{Pale}\texttt{auth\_sessions}\newline\texttt{user\_id} & \texttt{app\_users}\newline\texttt{id} & NO\newline\texttt{CASCADE} & 登录会话可以撤销，不能替代账户或患者主档。\\
\texttt{patient\_account\_links}\newline\texttt{patient\_id} & \texttt{patients}\newline\texttt{id} & NO\newline\texttt{CASCADE} & 患者主档与本人账户的绑定须经过身份核实。\\
\rowcolor{Pale}\texttt{patient\_account\_links}\newline\texttt{user\_id} & \texttt{app\_users}\newline\texttt{id} & NO\newline\texttt{CASCADE} & 账户与患者关系独立记录，便于处理解绑及复核。\\
\texttt{notification\_events}\newline\texttt{appointment\_id} & \texttt{appointments}\newline\texttt{id} & YES\newline\texttt{CASCADE} & 预约通知保留投递事实，发送失败不改变预约状态。\\
\bottomrule\end{longtable}\endgroup


\subsubsection{患者主档与风险资料：关键关系与约束}
患者全网主档是临床、预约与收费的共同锚点。紧急联系人、过敏和同意材料是可重复、可追溯的独立事实。

\begingroup\footnotesize\setlength{\tabcolsep}{3pt}
\begin{longtable}{@{}L{43mm}L{32mm}L{24mm}L{49mm}@{}}
\caption{患者主档与风险资料关系路径（Key Paths）}\label{tab:domain-key-03}\\\toprule
Child table / FK & Parent table / PK & Null; Delete & Business meaning / 业务含义\\\midrule\endfirsthead
\toprule Child table / FK & Parent table / PK & Null; Delete & Business meaning / 业务含义\\\midrule\endhead
\texttt{patients}\newline\texttt{primary\_clinic\_id} & \texttt{clinics}\newline\texttt{id} & YES\newline\texttt{SET NULL} & 首诊诊所为主档来源，不限制患者后续跨店就诊。\\
\rowcolor{Pale}\texttt{patient\_allergies}\newline\texttt{patient\_id} & \texttt{patients}\newline\texttt{id} & NO\newline\texttt{CASCADE} & 同一患者可记录多种过敏，复核记录不可覆盖原始事实。\\
\texttt{patient\_consents}\newline\texttt{patient\_id} & \texttt{patients}\newline\texttt{id} & NO\newline\texttt{CASCADE} & 每次同意材料与患者关联，并保留签署和适用范围。\\
\rowcolor{Pale}\texttt{emergency\_contacts}\newline\texttt{patient\_id} & \texttt{patients}\newline\texttt{id} & NO\newline\texttt{CASCADE} & 紧急联系人可重复，不能压成患者表的固定单字段。\\
\texttt{patient\_duplicate}\newline\texttt{reviews}\newline\texttt{source\_patient\_id} & \texttt{patients}\newline\texttt{id} & NO\newline\texttt{RESTRICT} & 疑似重复患者先进入复核流程，不自动合并主档。\\
\rowcolor{Pale}\texttt{patient\_timeline}\newline\texttt{events}\newline\texttt{clinic\_id} & \texttt{clinics}\newline\texttt{id} & YES\newline\texttt{SET NULL} & 时间线事件保留发生诊所，支持跨店追溯。\\
\bottomrule\end{longtable}\endgroup


\subsubsection{牙医、资质与治疗目录：关键关系与约束}
牙医任职、资质证书、治疗目录和项目所需资质分别建模；能否执业须同时考虑门店、项目和服务日期。

\begingroup\footnotesize\setlength{\tabcolsep}{3pt}
\begin{longtable}{@{}L{43mm}L{32mm}L{24mm}L{49mm}@{}}
\caption{牙医、资质与治疗目录关系路径（Key Paths）}\label{tab:domain-key-04}\\\toprule
Child table / FK & Parent table / PK & Null; Delete & Business meaning / 业务含义\\\midrule\endfirsthead
\toprule Child table / FK & Parent table / PK & Null; Delete & Business meaning / 业务含义\\\midrule\endhead
\texttt{dentist\_clinic}\newline\texttt{assignments}\newline\texttt{dentist\_id} & \texttt{dentists}\newline\texttt{id} & NO\newline\texttt{RESTRICT} & 任职关系独立于牙医主档，可表示跨店任职和有效期。\\
\rowcolor{Pale}\texttt{dentist\_clinic}\newline\texttt{assignments}\newline\texttt{clinic\_id} & \texttt{clinics}\newline\texttt{id} & NO\newline\texttt{RESTRICT} & 门店排班须核对该诊所任职关系。\\
\texttt{dentist}\newline\texttt{qualifications}\newline\texttt{qualification\_type\_id} & \texttt{qualification}\newline\texttt{types}\newline\texttt{id} & NO\newline\texttt{RESTRICT} & 证书类别与证书实例分开，便于到期和更新。\\
\rowcolor{Pale}\texttt{treatment}\newline\texttt{qualification}\newline\texttt{requirements}\newline\texttt{treatment\_id} & \texttt{treatment}\newline\texttt{catalog}\newline\texttt{id} & NO\newline\texttt{CASCADE} & 治疗项目与所需资质是多对多规则，不嵌入目录文本。\\
\texttt{treatment}\newline\texttt{qualification}\newline\texttt{requirements}\newline\texttt{qualification\_type\_id} & \texttt{qualification}\newline\texttt{types}\newline\texttt{id} & NO\newline\texttt{RESTRICT} & 同类资质要求可服务多个治疗项目。\\
\bottomrule\end{longtable}\endgroup


\subsubsection{排班、资源与预约：关键关系与约束}
可用时段由营业时间、停诊、牙医排班和牙椅占用共同决定；预约表只保存已确认的选择，冲突校验仍需事务执行。

\begingroup\footnotesize\setlength{\tabcolsep}{3pt}
\begin{longtable}{@{}L{43mm}L{32mm}L{24mm}L{49mm}@{}}
\caption{排班、资源与预约关系路径（Key Paths）}\label{tab:domain-key-05}\\\toprule
Child table / FK & Parent table / PK & Null; Delete & Business meaning / 业务含义\\\midrule\endfirsthead
\toprule Child table / FK & Parent table / PK & Null; Delete & Business meaning / 业务含义\\\midrule\endhead
\texttt{clinic\_chairs}\newline\texttt{clinic\_id} & \texttt{clinics}\newline\texttt{id} & NO\newline\texttt{RESTRICT} & 牙椅是诊所资源，分配前须核对所属门店。\\
\rowcolor{Pale}\texttt{dentist\_availability}\newline\texttt{rules}\newline\texttt{dentist\_id} & \texttt{dentists}\newline\texttt{id} & NO\newline\texttt{CASCADE} & 常规可用规则与单次预约分离。\\
\texttt{dentist\_schedule}\newline\texttt{exceptions}\newline\texttt{clinic\_id} & \texttt{clinics}\newline\texttt{id} & YES\newline\texttt{CASCADE} & 门店级排班例外可关闭某日资源。\\
\rowcolor{Pale}\texttt{appointments}\newline\texttt{patient\_id} & \texttt{patients}\newline\texttt{id} & NO\newline\texttt{RESTRICT} & 预约始终指向已识别患者，不能以姓名字符串代替关联。\\
\texttt{appointments}\newline\texttt{chair\_id} & \texttt{clinic\_chairs}\newline\texttt{id} & YES\newline\texttt{RESTRICT} & 牙椅为可选分配；空值不代表资源校验通过。\\
\rowcolor{Pale}\texttt{appointment\_status}\newline\texttt{history}\newline\texttt{appointment\_id} & \texttt{appointments}\newline\texttt{id} & NO\newline\texttt{CASCADE} & 状态变更追加留痕，不以当前状态覆盖历史。\\
\texttt{waitlist\_offers}\newline\texttt{waitlist\_entry\_id} & \texttt{waitlist\_entries}\newline\texttt{id} & NO\newline\texttt{CASCADE} & 候补邀约保留来源候补记录，便于追踪补位结果。\\
\bottomrule\end{longtable}\endgroup


\subsubsection{临床就诊与病历版本：关键关系与约束}
临床事实沿预约—就诊—实施项目展开。病历版本、过敏确认和附件独立保存，以区分草稿、签署和后续更正。

\begingroup\footnotesize\setlength{\tabcolsep}{3pt}
\begin{longtable}{@{}L{43mm}L{32mm}L{24mm}L{49mm}@{}}
\caption{临床就诊与病历版本关系路径（Key Paths）}\label{tab:domain-key-06}\\\toprule
Child table / FK & Parent table / PK & Null; Delete & Business meaning / 业务含义\\\midrule\endfirsthead
\toprule Child table / FK & Parent table / PK & Null; Delete & Business meaning / 业务含义\\\midrule\endhead
\texttt{clinical\_visits}\newline\texttt{appointment\_id} & \texttt{appointments}\newline\texttt{id} & NO\newline\texttt{RESTRICT} & 就诊回指预约，须另外使用唯一约束验证至多一次有效就诊。\\
\rowcolor{Pale}\texttt{clinical\_visits}\newline\texttt{patient\_id} & \texttt{patients}\newline\texttt{id} & NO\newline\texttt{RESTRICT} & 就诊患者应与预约患者一致，该跨表规则由事务核对。\\
\texttt{actual\_treatment}\newline\texttt{items}\newline\texttt{visit\_id} & \texttt{clinical\_visits}\newline\texttt{id} & NO\newline\texttt{CASCADE} & 一处就诊可产生多项实际治疗，不能只保留计划项目。\\
\rowcolor{Pale}\texttt{clinical\_note}\newline\texttt{versions}\newline\texttt{clinical\_note\_id} & \texttt{clinical\_notes}\newline\texttt{id} & NO\newline\texttt{CASCADE} & 病历每次修订形成版本，签署版本不原地覆盖。\\
\texttt{visit\_allergy}\newline\texttt{acknowledgements}\newline\texttt{allergy\_id} & \texttt{patient}\newline\texttt{allergies}\newline\texttt{id} & NO\newline\texttt{RESTRICT} & 到诊时对具体过敏记录作确认，不能只有笼统风险标志。\\
\bottomrule\end{longtable}\endgroup


\subsubsection{结算、收款与成本：关键关系与约束}
结算关系以实际治疗为来源。单据、明细、收款、退款和调整采用分表事实，净收款通过有效流水重算。

\begingroup\footnotesize\setlength{\tabcolsep}{3pt}
\begin{longtable}{@{}L{43mm}L{32mm}L{24mm}L{49mm}@{}}
\caption{结算、收款与成本关系路径（Key Paths）}\label{tab:domain-key-07}\\\toprule
Child table / FK & Parent table / PK & Null; Delete & Business meaning / 业务含义\\\midrule\endfirsthead
\toprule Child table / FK & Parent table / PK & Null; Delete & Business meaning / 业务含义\\\midrule\endhead
\texttt{invoices}\newline\texttt{visit\_id} & \texttt{clinical\_visits}\newline\texttt{id} & NO\newline\texttt{RESTRICT} & 一处就诊至多开具一张有效结算单，唯一性需落到约束。\\
\rowcolor{Pale}\texttt{invoice\_lines}\newline\texttt{actual\_treatment\_item}\newline\texttt{id} & \texttt{actual\_treatment}\newline\texttt{items}\newline\texttt{id} & NO\newline\texttt{RESTRICT} & 每条计费明细溯源实际治疗，防止仅凭目录价格开单。\\
\texttt{payments}\newline\texttt{invoice\_id} & \texttt{invoices}\newline\texttt{id} & NO\newline\texttt{RESTRICT} & 一张单据可有多笔收款，金额累计须核对余额。\\
\rowcolor{Pale}\texttt{payment\_refunds}\newline\texttt{payment\_id} & \texttt{payments}\newline\texttt{id} & NO\newline\texttt{RESTRICT} & 退款必须指向原收款，并校验累计可退余额。\\
\texttt{invoice\_adjustments}\newline\texttt{approved\_by\_user\_id} & \texttt{app\_users}\newline\texttt{id} & YES\newline\texttt{SET NULL} & 授权调整保留审批人；未批准不能进入正式核算。\\
\rowcolor{Pale}\texttt{treatment\_cost}\newline\texttt{entries}\newline\texttt{actual\_treatment\_item}\newline\texttt{id} & \texttt{actual\_treatment}\newline\texttt{items}\newline\texttt{id} & NO\newline\texttt{RESTRICT} & 用品成本与实施项目关联，不能由账单金额倒推。\\
\bottomrule\end{longtable}\endgroup


\subsubsection{期间核算与激励：关键关系与约束}
期间、门店和网络三个核算层级须分开。期间核定结果和激励池是快照，缺失成本或费用凭据时先生成关闭异常。

\begingroup\footnotesize\setlength{\tabcolsep}{3pt}
\begin{longtable}{@{}L{43mm}L{32mm}L{24mm}L{49mm}@{}}
\caption{期间核算与激励关系路径（Key Paths）}\label{tab:domain-key-08}\\\toprule
Child table / FK & Parent table / PK & Null; Delete & Business meaning / 业务含义\\\midrule\endfirsthead
\toprule Child table / FK & Parent table / PK & Null; Delete & Business meaning / 业务含义\\\midrule\endhead
\texttt{operating\_expenses}\newline\texttt{reporting\_period\_id} & \texttt{financial}\newline\texttt{reporting}\newline\texttt{periods}\newline\texttt{id} & NO\newline\texttt{RESTRICT} & 费用归属明确核算期间，不能只按录入日统计。\\
\rowcolor{Pale}\texttt{financial\_close}\newline\texttt{exceptions}\newline\texttt{reporting\_period\_id} & \texttt{financial}\newline\texttt{reporting}\newline\texttt{periods}\newline\texttt{id} & NO\newline\texttt{RESTRICT} & 关闭异常指向待核定期间，异常未消除不得伪造完整结果。\\
\texttt{financial\_period}\newline\texttt{results}\newline\texttt{clinic\_id} & \texttt{clinics}\newline\texttt{id} & YES\newline\texttt{RESTRICT} & 门店结果与网络结果的范围必须明确区分。\\
\rowcolor{Pale}\texttt{incentive\_pools}\newline\texttt{network\_result\_id} & \texttt{financial\_period}\newline\texttt{results}\newline\texttt{id} & NO\newline\texttt{RESTRICT} & 激励池依据已核定网络结果，而非实时估算。\\
\texttt{incentive\_pools}\newline\texttt{approved\_by\_user\_id} & \texttt{app\_users}\newline\texttt{id} & YES\newline\texttt{SET NULL} & 批准责任独立留存，计算结果不自动等于可发放结果。\\
\bottomrule\end{longtable}\endgroup


\subsubsection{扩张情景与候选诊所：关键关系与约束}
扩张测算是情景分析，不是已批准开店事实。假设、候选诊所和模拟结果分别保存，使每次预测可复算。

\begingroup\footnotesize\setlength{\tabcolsep}{3pt}
\begin{longtable}{@{}L{43mm}L{32mm}L{24mm}L{49mm}@{}}
\caption{扩张情景与候选诊所关系路径（Key Paths）}\label{tab:domain-key-09}\\\toprule
Child table / FK & Parent table / PK & Null; Delete & Business meaning / 业务含义\\\midrule\endfirsthead
\toprule Child table / FK & Parent table / PK & Null; Delete & Business meaning / 业务含义\\\midrule\endhead
\texttt{expansion\_scenarios}\newline\texttt{created\_by\_user\_id} & \texttt{app\_users}\newline\texttt{id} & YES\newline\texttt{SET NULL} & 情景创建人用于方案责任追溯。\\
\rowcolor{Pale}\texttt{expansion\_scenarios}\newline\texttt{approved\_by\_user\_id} & \texttt{app\_users}\newline\texttt{id} & YES\newline\texttt{SET NULL} & 审批人可选，未批准的情景不得标成已定方案。\\
\texttt{expansion\_assumptions}\newline\texttt{scenario\_id} & \texttt{expansion}\newline\texttt{scenarios}\newline\texttt{id} & NO\newline\texttt{CASCADE} & 测算参数归属特定情景，不覆盖别的版本。\\
\rowcolor{Pale}\texttt{expansion\_candidate}\newline\texttt{clinics}\newline\texttt{scenario\_id} & \texttt{expansion}\newline\texttt{scenarios}\newline\texttt{id} & NO\newline\texttt{CASCADE} & 一个情景可比较多个候选地点。\\
\texttt{expansion\_scenario}\newline\texttt{results}\newline\texttt{scenario\_id} & \texttt{expansion}\newline\texttt{scenarios}\newline\texttt{id} & NO\newline\texttt{CASCADE} & 预测结果与输入情景关联，展示时注明非实际经营数据。\\
\bottomrule\end{longtable}\endgroup

\noindent\emph{说明：}上表精选的是能够解释九张分域图主线的关系，不以节省篇幅隐藏其余外键。
附属关系清单记录现有 119 条外键的完整目标、字段和删除动作；
图中未显示的三个目标增强外键仍按本节后文的迁移建议处理，未混入现有关系。
